\subsection{Closing the rank-one case for every bounded scalar complex}

The remaining nonflat-cohomology alternative in the scalar
replacement can be excluded in the rank-one product by an
elementary two-variable argument. This proof is retained
independently of the subsequent higher-dimensional extension.

Let \(k\) have characteristic zero and let
\[
 A=k[x,y],\qquad \deg x=\deg y=2,\qquad
 W=\langle s_1,s_2\rangle\simeq C_2\times C_2,
\]
where \(s_1\) negates \(x\), \(s_2\) negates \(y\), and
\(\tau=s_1s_2\) negates both. The ordinary Soergel generators are
\[
 B_i=A\otimes_{A^{s_i}}A(1),\qquad b_i=[B_i].
\]
Retain \(\delta=v+v^{-1}\).

\begin{finding}[\fid{F59}, arbitrary bounded scalar coefficients in rank one]
\label{f:rankOneScalarClosure}
There is no bounded ordinary Soergel complex \(X\) with
termwise split-Grothendieck Euler class
\begin{equation}\label{eq:rankOneScalarClass}
 [X]=b_1b_2+(b_1-\delta)(b_2-\delta),
\end{equation}
and no bounded complex \(C\) of finite graded \(k\)-vector
spaces with
\begin{equation}\label{eq:rankOneScalarEuler}
 \chi_v(C)=\sum_{j,e}(-1)^j\dim_kH^j(C)_e\,v^e
         =\delta^2,
\end{equation}
admitting an actual derived bimodule isomorphism
\begin{equation}\label{eq:rankOneScalarQuadratic}
 X\otimes_A^{\mathbf L}X\simeq C\otimes_kX.
\end{equation}
In particular, no homotopy quadratic with these objects
exists. Arbitrary bounded amplitude and nonflat
cohomology do not repair this ordinary rank-one lift.
\end{finding}

\begin{proof}
\par\smallskip\noindent\textbf{Generic positive cohomology.}\quad
Put \(Q=\operatorname{Frac}(A)\) and forget internal
grading at this stage. Ordinary exact monoidal graph
localization has simples \(Q_g\), \(g\in W\), with
\[
 Q_g\otimes_QQ_h\simeq Q_{gh},\qquad
 \operatorname{Hom}(Q_g,Q_h)=0\quad(g\ne h).
\]
The exact primary framework is
\cite[\S1.1, equations (1.1)--(1.2);
\S1.2, Theorem~1.6 and Remark~1.7]{EWLocalized}.
All terms of \(X\) are free on both sides.
Let
\[
 F(u)=\sum_{g\in W}f_g(u)g,\quad
 f_g(u)\in\mathbb N[u,u^{-1}],\qquad
 d(u)=\sum_j c_j u^j,\quad c_j=\dim_kH^j(C),
\]
where \(u\) records actual cohomological degree.
Generic K\"unneth and the actual quadratic give
\begin{equation}\label{eq:rankOneGenericQuadratic}
 F(u)^2=d(u)F(u).
\end{equation}
The two Euler classes imply
\begin{equation}\label{eq:rankOneGenericEuler}
 F(-1)=2e+2\tau,\qquad d(-1)=4.
\end{equation}
Indeed \(b_i\) gives \(e+s_i\) and
\(b_i-\delta\) gives \(s_i-e\) after forgetting shifts.
In particular both polynomials are nonzero.

For completeness, augmenting
\eqref{eq:rankOneGenericQuadratic} gives \(a^2=da\),
where \(a=\sum_gf_g\); hence \(a=d\).
For every real \(u>0\), the numbers
\(\mu_u(g)=f_g(u)/d(u)\) form an idempotent probability
measure. Its support \(S\) is independent of \(u\).
Positivity gives \(S^2=S\), so this nonempty finite
semigroup inside a group is a subgroup. If \(g_0\)
maximizes \(\mu(g)\), then
\[
 \mu(g_0)=\sum_{x\in S}\mu(x)\mu(x^{-1}g_0)
\]
is a weighted average of quantities at most that maximum,
with all weights positive. Thus every value on \(S\)
is the same. The measure is uniform.
Equality for all \(u>0\) is a Laurent-polynomial identity.
At \(u=-1\), \eqref{eq:rankOneGenericEuler} forces
\(|S|=2\) and \(S=\{e,\tau\}\). Therefore
\begin{equation}\label{eq:rankOneUniformGraphs}
 F(u)=\frac{d(u)}2(e+\tau),\qquad
 \operatorname{rank}_A M_j=c_j,\quad M_j=H^j(X).
\end{equation}
Every generic cohomology module has support contained
in these two graphs; its rank may still be zero.

This conclusion has deliberately forgotten internal
grading. A degree-two homogeneous unit after localization
identifies shifts differing by two. Full coefficientwise
parity in \(v\) cannot be inferred here; it will be
recovered only from global free graph summands.

\par\smallskip\noindent\textbf{The two-variable tensor-dimension lemma.}\quad
Let \(\mathfrak m=(x,y)\). If \(M\) is a finite graded
left \(A\)-module of rank \(r\), and \(V\ne0\) is a
finite-dimensional graded right \(A\)-module with its
positive-degree actions, then
\begin{equation}\label{eq:rankOneTensorDimension}
 \dim_k(V\otimes_AM)\ge r\dim_kV,
\end{equation}
with equality exactly when \(M\) is graded free,
including the zero module.

Here is a complete proof. First suppose \(M\) is
torsion-free. Construct its minimal graded free
resolution by taking homogeneous bases modulo
\(\mathfrak m\) for successive syzygies. Each differential
has entries in \(\mathfrak m\), so
\[
 F_i/\mathfrak mF_i\simeq\operatorname{Tor}_i^A(k,M).
\]
The explicit Koszul resolution
\[
 0\longrightarrow A
 \xrightarrow{\,a\mapsto(-ya,xa)\,}A^2
 \xrightarrow{\,(b,c)\mapsto xb+yc\,}A
 \longrightarrow k\longrightarrow0
\]
has its corresponding internal shifts and is exact
by coprimality. Tensoring by \(M\) gives vanishing
Tor in degrees above two, while
\(\operatorname{Tor}_2(k,M)=\{m:xm=ym=0\}\), up to
shift, vanishes by torsion-freeness.
Minimality and graded Nakayama give
\[
 0\longrightarrow F_1\longrightarrow F_0
       \longrightarrow M\longrightarrow0.
\]
Writing \(f_i=\operatorname{rank}F_i\), localization
gives \(r=f_0-f_1\), and tensoring with \(V\) gives
\begin{equation}\label{eq:rankOneTorDimension}
 \dim(V\otimes_AM)
 =r\dim V+\dim\operatorname{Tor}_1^A(V,M).
\end{equation}
If \(M\) is nonfree, \(F_1\ne0\).
A nonzero vector of highest degree in \(V\) lies in
\(\operatorname{soc}(V)=\{v:vx=vy=0\}\).
All differential entries belong to \(\mathfrak m\),
so a copy of this socle for each generator of \(F_1\)
lies in the tensor differential's kernel.
Thus \(\operatorname{Tor}_1(V,M)\ne0\), proving
strictness. A free module gives equality.

For arbitrary \(M\), let \(T\) be its full torsion
submodule and \(Q_0=M/T\). The submodule is graded:
inspect the highest-degree equation of a polynomial
annihilator, remove that torsion component, and descend.
The quotient is torsion-free, since if \(az\) is
torsion for nonzero \(a\), then \(baz=0\) for some
nonzero \(b\). Both are finite graded modules and
\(\operatorname{rank}Q_0=r\).
If \(Q_0\) is nonfree, strictness for \(Q_0\) and
the surjection \(V\otimes M\to V\otimes Q_0\) prove
strictness for \(M\).
If \(Q_0\) is free, lifting a homogeneous basis
splits \(M\to Q_0\) as a left-module map.
When \(T\ne0\), the tensor \(V\otimes T\) surjects onto
\[
 (V/V\mathfrak m)\otimes_k(T/\mathfrak mT)\ne0.
\]
Thus strictness holds again. If \(T=0\), \(M\) is
free. This proves \eqref{eq:rankOneTensorDimension}
in all cases, using only a left-module splitting.

For every nonzero finite graded bimodule \(M\) that
is finite on the left,
\begin{equation}\label{eq:rankOneModuleSquareNonzero}
                         M\otimes_AM\ne0.
\end{equation}
Indeed \(V=k\otimes_AM\), using the left action,
is a nonzero finite graded right origin module.
The left-origin fiber of the square is \(V\otimes M\),
which surjects onto
\((V/V\mathfrak m)\otimes_k(M/\mathfrak mM)\ne0\).
The positive grading and common origin are essential.
Moreover, if
\[
 M\otimes_AM\simeq U\otimes_kM,\qquad
 \dim_kU=\operatorname{rank}_AM,
\]
then taking this origin fiber gives equality in
\eqref{eq:rankOneTensorDimension}, so \(M\) is left-free.

\par\smallskip\noindent\textbf{Global graph splitting.}\quad
A finite graded left-free bimodule \(M\) with generic
support only \(e,\tau\) is a sum of actual shifted
\(A_e,A_\tau\). To prove this, let \(U,V\) be the
right multiplication matrices of \(x,y\). Generically,
\[
 yU=xV,\qquad U^2=x^2I.
\]
Left torsion-freeness extends the identities to \(A\).
Coprimality gives \(U=xZ,V=yZ\), with \(Z\) homogeneous
of degree zero and \(Z^2=I\).
The bimodule idempotents \((I+Z)/2,(I-Z)/2\) have
respectively the ordinary and \(\tau\)-twisted actions.
Their images are finite graded left-projective.
Such a module is graded free: lift a homogeneous
residue basis, split the resulting free surjection,
and apply graded Nakayama to its complementary kernel.
The two images are therefore sums of graph modules
and are also right-free.

\par\smallskip\noindent\textbf{The global top is visible and free.}\quad
Let \(a=\max H(X)\), \(c=\max H(C)\); the nonzero
Euler classes make these finite and defined.
The top-truncation identity gives
\[
 H^{2a}(X\otimes_A^{\mathbf L}X)
              \simeq M_a\otimes_AM_a\ne0.
\]
There is no higher cohomology. This is the
top-cokernel identity for bounded-above flat
resolutions, not a general K\"unneth splitting.
The scalar right side of
\eqref{eq:rankOneScalarQuadratic} has top degree
\(a+c\), so \(c=a\).
Let \(b\) be the highest generic cohomology degree.
Its nonzero graph top has nonzero square, and the
generic quadratic gives \(2b=b+c\).
Thus \(b=c=a\). In particular
\[
 r=\operatorname{rank}M_a=c_a>0,
 \qquad
 M_a\otimes_AM_a\simeq H^a(C)\otimes_kM_a.
\]
The tensor-dimension equality proves that \(M_a\)
is left-free, and global graph splitting makes it
right-free as well.
Global nonvanishing was proved before its survival
under localization was used.

\par\smallskip\noindent\textbf{The canonical near-top K\"unneth map.}\quad
Suppose \(M_i\) is right-projective for every \(i>s\),
with \(s<a\). Then the canonical cycle cross-product
\begin{equation}\label{eq:rankOneNearTopKunneth}
 \bigoplus_{i+j=a+s}M_i\otimes_AM_j
       \longrightarrow H^{a+s}(X\otimes_A^{\mathbf L}X)
\end{equation}
is a graded bimodule isomorphism.

To verify this without bimodule formality, forget the
left action. As a bounded complex of projective right
modules, \(X\) can have each higher projective cohomology
split off starting at the top. Remove acyclic top
projective tails, choose sections onto the top projective
cohomology, cancel the resulting contractible pairs,
and repeat. This gives a right-module homotopy decomposition
\[
 X\simeq\left(\bigoplus_{i>s}M_i[-i]\right)\oplus X_{\le s},
\]
where the last right-projective complex has cohomology
at most \(s\) and top cohomology \(M_s\).
Tensor with the second \(X\). The higher right-flat
summands contribute \(M_i\otimes M_{a+s-i}\).
The remainder contributes \(M_s\otimes M_a\) by
the top-truncation identity. These are all terms
in \eqref{eq:rankOneNearTopKunneth}.
Their maps are the ordinary cycle cross-product.
The auxiliary one-sided splitting proves that this
canonical map is bijective; the map itself is
bimodule-linear. Hence the resulting isomorphism
is a bimodule statement. No bimodule splitting of
the complex is asserted.

\par\smallskip\noindent\textbf{Descending freeness.}\quad
For a finite graded left module \(N\), put
\[
 \phi(N)=\dim_k(N/\mathfrak mN).
\]
This is finite, detects zero by graded Nakayama,
and equals rank on a free module.
Assume all \(M_i\), \(i>s\), are free graph modules,
as established initially for \(M_a\).
Put \(N=M_s\), \(z=\phi(N)\), and retain \(r=c_a>0\).
The two terms involving \(N\) in
\eqref{eq:rankOneNearTopKunneth} are
\(M_a\otimes N\) and \(N\otimes M_a\).
Each has origin-fiber dimension \(rz\).
Right tensoring by a graph twists only the right
action; left tensoring twists the left action by
\(e\) or \(\tau\), which preserves \(\mathfrak m\).
All other terms are products of known free graphs
and contribute
\[
 S_s=\sum_{\substack{i+j=a+s\\i>s,\ j>s}}c_i c_j.
\]
Thus the left side of the actual cohomology
isomorphism has fiber dimension \(2rz+S_s\).

Scalar K\"unneth on the right gives
\[
 H^{a+s}(C\otimes_kX)
       =\bigoplus_{j+i=a+s}H^j(C)\otimes_kM_i.
\]
The terms \(j=a,i=s\) and \(j=s,i=a\) contribute
\(rz\) and \(rc_s\), respectively.
Every other contribution is \(S_s\), since
\(\operatorname{rank}M_i=c_i\).
Applying the ordinary left-origin fiber to the
actual isomorphism gives
\begin{equation}\label{eq:rankOneOriginInduction}
 2rz+S_s=rz+rc_s+S_s,\qquad
 \phi(M_s)=c_s=\operatorname{rank}M_s.
\end{equation}
If \(c_s=0\), graded Nakayama gives \(M_s=0\).
Otherwise a minimal graded free surjection onto
\(M_s\) has rank \(c_s\). Its kernel has rank zero
inside a free module over a domain, and is therefore
zero. Thus \(M_s\) is left-free.
Global graph splitting also makes it right-free,
allowing the induction to continue.
All cohomology modules are consequently sums of
actual shifted \(A_e,A_\tau\).

\par\smallskip\noindent\textbf{Restoring global internal grading.}\quad
There are now actual nonnegative Laurent polynomials
\[
 a(v,u),b(v,u)\in\mathbb N[v,v^{-1},u,u^{-1}]
\]
such that global cohomology has character \(ae+b\tau\),
with \(b\ne0\). Put
\[
 d(v,u)=\sum_{j,e}\dim_kH^j(C)_e\,v^eu^j.
\]
All cohomology is right-projective, so the preceding
one-sided splitting proves the canonical cross-product
K\"unneth map is a bimodule isomorphism in every degree.
Different graph types have no nonzero maps: a coefficient
of such a map is killed by \(2x\) in a domain.
Within each graph type, the graded origin fiber
determines the multiplicities. The actual quadratic gives
\[
 a^2+b^2=da,\qquad 2ab=db.
\]
The integral Laurent ring is a domain and \(b\ne0\),
so \(d=2a\). At \(u=-1\),
\[
                    v^2+2+v^{-2}=2a(v,-1),
\]
which is impossible because the two outer coefficients
on the left are odd.
The full internal grading has been used only after
global free graph decompositions were proved.
\end{proof}

\paragraph{Exact scope.}
The coefficient is genuinely scalar over \(k\).
A unit-bimodule complex with polynomial differentials,
a localization removing the origin, an unbounded
completion, or a changed category is outside the theorem.
A uniform all-ranks ordinary construction must include
this case and is excluded already here.
This elementary proof alone does not claim that every
isolated higher-rank construction has a compatible
rank-one restriction.
With extra fixed central variables in a rank-one
permutation realization, an actual homotopy quadratic
descends by termwise quotient: the objects, inverse
maps and homotopies all descend. No claim that an
arbitrary underived quotient preserves every derived
isomorphism is used. No arithmetic purity or
canonical-basis identification enters the argument.

% Separate paragraph for Part III.
\paragraph{The scalar-complex closure in Grothendieck's framework.}
The affine-plane module/sheaf dictionary is EGA~I,
Theorem~1.3.7, p.~86, Corollary~1.3.8, p.~87,
and Theorem~1.4.1, pp.~90--92.
The bounded free complexes fit SGA~6, I,
Definition~4.7, pp.~121--122 and Example~4.8, p.~122
\cite{EGA,SGA6}. The actual origin fiber, Tor groups
and derived quadratic retain information that an
Euler class discards; the proof establishes cohomological
freeness before using it.

The internal grading is a split-\(\mathbb G_m\) action
in SGA~3, I, 4.7.2--3, with cote~90, scan~17.
The actual effective coefficient classes and their
passage to virtual \(K_0\) belong to the representation
framework of SGA~6, V, \S\S2--3, with cote~133, scan~18
\cite{SGA3,SGA6,Archive}. These anchors do not assert
\(\gamma\)-filtration positivity or a \(\lambda\)-ring
on the noncommutative word algebra.
The ordinary Soergel normalization is
\cite[\S3.3, p.~1105; \S3.4, pp.~1106--1107;
\S3.5, Theorem~3.5]{EW}.
The tensor-dimension, canonical near-top K\"unneth,
descending-freeness and parity arguments are proved
here; they are not attributed to an unpublished manuscript.
